\documentclass[11pt,a4paper]{article}
\usepackage[utf8]{inputenc}
\usepackage[T1]{fontenc}
\usepackage[french]{babel}
\usepackage{amsmath, amssymb, amsthm}
\usepackage{geometry}
\geometry{margin=2.5cm}
\usepackage{hyperref}

\newcommand{\Omega}{\Omega}
\newcommand{\dd}{\,\mathrm{d}}

\title{\textbf{Projet VMAM 2025--2026}\\Réponses théoriques}
\author{Casals Diego}
\date{\today}

\begin{document}
\maketitle
\tableofcontents
\bigskip

% ============================================================
\section{Problème direct}

\subsection{Formulation variationnelle}

On cherche $T \in V$ où $V = \{v \in H^1(\Omega) \mid v|_{\Gamma_{\text{haut}}} = T_c\}$.
Les conditions aux limites sont :
\begin{itemize}
  \item Bords gauche et droit ($\Gamma_{\text{lat}}$) : condition de Neumann homogène $\partial_n T = 0$ (parois isolées).
  \item Bord supérieur ($\Gamma_{\text{haut}}$) : condition de Dirichlet $T = T_c = 343\,\mathrm{K}$.
  \item Bord inférieur ($\Gamma_{\text{bas}}$, porte) : condition de Robin $\partial_n T + h(T - T_e) = 0$.
\end{itemize}

En multipliant l'équation $-\mathrm{div}(k\nabla T) = f$ par une fonction test $v \in V_0 = \{v \in H^1(\Omega) \mid v|_{\Gamma_{\text{haut}}} = 0\}$ et en intégrant par parties, on obtient :
\[
\int_\Omega k\,\nabla T \cdot \nabla v \dd x - \int_{\partial\Omega} k\,\partial_n T\, v \dd\sigma = \int_\Omega f\,v \dd x.
\]
En substituant les conditions aux limites (Neumann nul sur $\Gamma_{\text{lat}}$, Robin sur $\Gamma_{\text{bas}}$, et $v=0$ sur $\Gamma_{\text{haut}}$) :

\begin{equation}\label{eq:VF}
\boxed{
\int_\Omega k\,\nabla T \cdot \nabla v \dd x + \int_{\Gamma_{\text{bas}}} h\,T\,v \dd\sigma
= \int_\Omega f\,v \dd x + \int_{\Gamma_{\text{bas}}} h\,T_e\,v \dd\sigma,
\quad \forall v \in V_0.
}
\end{equation}

\paragraph{Analyse.}
La forme bilinéaire $a(T,v) = \int_\Omega k\nabla T\cdot\nabla v\dd x + \int_{\Gamma_{\text{bas}}} hTv\dd\sigma$ est :
\begin{itemize}
  \item \textit{Continue} sur $H^1(\Omega)$ par les inégalités de Cauchy-Schwarz et de trace.
  \item \textit{Coercive} sur $V_0$ : car $k \geq 1 > 0$ et $h > 0$, la condition de Dirichlet sur $\Gamma_{\text{haut}}$ assure que la semi-norme $H^1$ est une norme équivalente à la norme $H^1$ sur $V_0$ (inégalité de Poincaré).
\end{itemize}
Le membre de droite est une forme linéaire continue sur $V_0$. Le théorème de Lax-Milgram garantit donc l'existence et l'unicité de la solution.

\subsection{Approximation par éléments finis et système linéaire}

On approche $V_0$ par un espace de dimension finie $V_h \subset V_0$ basé sur un maillage triangulaire $\mathcal{T}_h$ de $\Omega$ et des éléments finis $\mathbb{P}_1$ (fonctions continues affines par morceaux). Soit $(\varphi_j)_{j=1}^N$ la base nodale de $V_h$.

On écrit $T_h = \sum_{j=1}^N T_j\,\varphi_j$ et on cherche les $T_j$ tels que~\eqref{eq:VF} soit satisfait pour tout $v = \varphi_i$. Cela conduit au système linéaire :
\[
\mathbf{A}\,\mathbf{T} = \mathbf{b},
\]
où :
\[
A_{ij} = \int_\Omega k\,\nabla\varphi_j\cdot\nabla\varphi_i\dd x + \int_{\Gamma_{\text{bas}}} h\,\varphi_j\,\varphi_i\dd\sigma,
\qquad
b_i = \int_\Omega f\,\varphi_i\dd x + \int_{\Gamma_{\text{bas}}} h\,T_e\,\varphi_i\dd\sigma.
\]
La condition de Dirichlet $T = T_c$ sur $\Gamma_{\text{haut}}$ est imposée en modifiant les lignes correspondantes de $\mathbf{A}$ et $\mathbf{b}$ (élimination des degrés de liberté Dirichlet). La matrice $\mathbf{A}$ est symétrique définie positive, creuse, et le système est résolu par un solveur direct (LU).

\subsection{Implémentation FreeFem++ et commentaires}

L'implémentation utilise FreeFem++ avec des éléments $\mathbb{P}_1$. La géométrie est construite via les commandes \texttt{border} et \texttt{buildmesh}, avec un raffinement autour des résistances et de la zone $S$.

Le coefficient de diffusion est défini par \texttt{kk = 1 + 9*(region != regS)}, soit $k=1$ dans $S$ et $k=10$ ailleurs. La condition de Robin sur la porte (label 1) est traitée par le terme \texttt{int1d(Th,1)(hc*uh*vh - hc*Te*vh)}, et la condition de Dirichlet $T=T_c$ sur le haut (label 3) par \texttt{on(3, uh=Tc)}.

Pour $\alpha_i = 10^4$ (i=1,\ldots,6), on observe une température plus élevée autour des résistances et une distribution quasi-uniforme dans $S$. On peut vérifier que la température augmente bien depuis la porte (condition de refroidissement) vers le haut (paroi chaude).

\subsection{Problème transitoire (phase de chauffe)}

Pour décrire la montée en température, on ajoute le terme instationnaire $\rho c\,\partial_t T$, ce qui donne :
\[
\rho c\,\partial_t T - \mathrm{div}(k\nabla T) = f \quad \text{dans } \Omega \times (0,t_{\max}).
\]
On approche la dérivée temporelle par un schéma d'Euler implicite de pas $\Delta t$ : $\partial_t T^{n+1} \approx (T^{n+1} - T^n)/\Delta t$. La formulation variationnelle à chaque pas de temps est :
\[
\int_\Omega \frac{\rho c}{\Delta t}T^{n+1}v\dd x + \int_\Omega k\nabla T^{n+1}\cdot\nabla v\dd x + \int_{\Gamma_{\text{bas}}} hT^{n+1}v\dd\sigma
= \int_\Omega \left(\frac{\rho c}{\Delta t}T^n + f\right)v\dd x + \int_{\Gamma_{\text{bas}}} hT_e\,v\dd\sigma.
\]
La condition initiale est $T^0 = T_e = 293\,\mathrm{K}$ (four initialement à température ambiante). Le schéma est inconditionnellement stable (Euler implicite), et la solution converge vers la solution stationnaire quand $t\to+\infty$.

% ============================================================
\section{Problème d'optimisation}

\subsection{Décomposition linéaire de la solution}

Par linéarité de l'équation~\eqref{eq:VF}, on décompose $f = \sum_{i=1}^6 \alpha_i \mathbf{1}_{C_i}$ et on écrit :
\[
T(\alpha, x) = T_0(x) + \sum_{i=1}^6 \alpha_i\,T_i(x),
\]
où $T_0$ est la solution du problème \emph{sans source} ($f=0$) avec les mêmes conditions aux limites (Dirichlet et Robin), et pour $i=1,\ldots,6$, $T_i$ est la solution du problème homogène en température de Dirichlet ($T_i|_{\Gamma_{\text{haut}}}=0$) avec le second membre $f = \mathbf{1}_{C_i}$ :
\[
-\mathrm{div}(k\nabla T_i) = \mathbf{1}_{C_i} \quad \text{dans } \Omega, \quad
T_i|_{\Gamma_{\text{haut}}} = 0, \quad
\partial_n T_i + h T_i = 0 \text{ sur } \Gamma_{\text{bas}}, \quad
\partial_n T_i = 0 \text{ sur } \Gamma_{\text{lat}}.
\]
Ces 7 problèmes sont indépendants des coefficients $(\alpha_i)$ et n'ont besoin d'être résolus qu'une seule fois.

\subsection{Réduction à un système linéaire}

On cherche $\alpha^* = (\alpha_i^*)_{i=1}^6$ minimisant la fonctionnelle :
\[
J(\alpha) = \frac{1}{2}\int_S |T(\alpha,x) - T_s|^2\dd x.
\]
En substituant la décomposition, $J$ devient une fonction quadratique convexe en $\alpha$ :
\[
J(\alpha) = \frac{1}{2}\int_S \left|T_0 + \sum_{i=1}^6 \alpha_i T_i - T_s\right|^2\dd x.
\]
Les conditions d'optimalité $\partial J/\partial \alpha_j = 0$ (pour $j=1,\ldots,6$) donnent :
\[
\int_S \left(T_0 + \sum_{i=1}^6 \alpha_i T_i - T_s\right)T_j\dd x = 0, \quad j = 1,\ldots,6,
\]
soit le système linéaire $6\times 6$ :
\begin{equation}\label{eq:syst_opt}
\boxed{M\,\alpha^* = b, \qquad M_{ij} = \int_S T_i\,T_j\dd x, \qquad b_j = \int_S (T_s - T_0)\,T_j\dd x.}
\end{equation}
La matrice $M$ est symétrique semi-définie positive. Elle est définie positive si les fonctions $(T_i|_S)_{i=1}^6$ sont linéairement indépendantes, ce qui est le cas générique.

\subsection{Coefficients optimaux pour $T_s = 400\,\mathrm{K}$}

Le système~\eqref{eq:syst_opt} est résolu par élimination de Gauss avec pivot partiel (implémentée manuellement dans le code FreeFem++). La solution $T^{\text{opt}} = T_0 + \sum_i \alpha_i^* T_i$ est ensuite comparée à une simulation directe utilisant ces coefficients, ce qui valide la procédure. La qualité est mesurée par l'écart-type $\sigma_S = \left(\frac{1}{|S|}\int_S (T^{\text{opt}} - T_s)^2\dd x\right)^{1/2}$ et par la valeur de $J(\alpha^*)$.

\subsection{Amélioration de la distribution par ajout de sources}

En augmentant le nombre de résistances (passage à 8 sources dans le code, placées sur une grille $\{-1, -1/3, 1/3, 1\} \times \{-0.75, 0.75\}$), le même formalisme linéaire s'applique avec un système $8\times8$. En jouant sur la disposition et la forme des sources, on peut réduire significativement l'écart-type de température dans $S$, au prix d'une légère augmentation du coût (intensités plus élevées). Des sources rapprochées du centre de $S$ sont plus efficaces, mais les résistances étant placées à l'extérieur de $S$, un compromis géométrique est nécessaire.

% ============================================================
\section{Problème inverse}

\subsection{Cadre et approche numérique}

On suppose qu'une source unique $f = \alpha_1 \mathbf{1}_{C_1}$ est inconnue (intensité $\alpha_1$ et localisation $C_1$ inconnues). On dispose de mesures de $T$ sur les bords latéraux $\Gamma_{\text{lat}}$ (ou de $\partial_n T$ sur $\Gamma_{\text{haut}}$).

On modélise la source par une gaussienne de position $(x_c, y_c)$ et d'écart-type $\sigma = r_0$ (pour s'affranchir de la non-différentiabilité de $\mathbf{1}_{C_1}$) :
\[
f(x, y; x_c, y_c, \alpha) = \alpha\,\exp\!\left(-\frac{(x-x_c)^2+(y-y_c)^2}{2\sigma^2}\right).
\]

\paragraph{Fonctionnelle de moindres carrés.}
Pour une position candidate $(x_c, y_c)$ fixée, on cherche $\alpha$ minimisant l'écart aux mesures sur $\Gamma_{\text{lat}}$ :
\[
\mathcal{J}(x_c, y_c, \alpha) = \frac{1}{2}\int_{\Gamma_{\text{lat}}} \left|T(x_c,y_c,\alpha) - T^{\text{mes}}\right|^2\dd\sigma.
\]
Par linéarité, $T = T_0 + \alpha\,\tilde{T}$ où $\tilde{T}$ est la solution unitaire ($\alpha=1$) associée à la source gaussienne. L'intensité optimale pour une position donnée est donc :
\[
\alpha^*(x_c,y_c) = \frac{\displaystyle\int_{\Gamma_{\text{lat}}} (T^{\text{mes}}-T_0)\,\tilde{T}\dd\sigma}{\displaystyle\int_{\Gamma_{\text{lat}}} \tilde{T}^2\dd\sigma}.
\]

\paragraph{Balayage sur grille.}
La position $(x_c, y_c)$ étant le paramètre non-linéaire, on procède par balayage sur une grille de $15\times12$ points couvrant $[-1.3, 1.3]\times[-0.8, 0.8]$. Pour chaque point de la grille, on résout le problème direct ($\tilde{T}$), on calcule $\alpha^*$ par la formule ci-dessus, puis on évalue $\mathcal{J}$. La position retenue est celle qui minimise $\mathcal{J}$.

\paragraph{Résultat.}
La position et l'intensité retrouvées sont proches des valeurs vraies (erreur de position de l'ordre du pas de grille, erreur d'intensité de quelques pourcents), ce qui valide l'approche. Des mesures aux deux bords latéraux (labels 2 et 4 dans FreeFem++) assurent une meilleure observabilité qu'un seul bord.

\end{document}
